\documentclass[11pt,letterpaper]{article}

% ── Geometry ──────────────────────────────────────────────────────────────────
\usepackage[margin=1in]{geometry}

% ── Pagination ────────────────────────────────────────────────────────────────
\usepackage[all]{nowidow}    % Require ≥2 lines on each side of a page break
\brokenpenalty=10000         % No hyphenated words across page breaks
\raggedbottom                % Allow uneven page bottoms rather than stretching

% ── Fonts ─────────────────────────────────────────────────────────────────────
\usepackage[T1]{fontenc}
\usepackage{newpxtext}     % Palatino-based serif (TeX Gyre Pagella)
\usepackage{newpxmath}     % Matching math font

% ── Colors ────────────────────────────────────────────────────────────────────
\usepackage[dvipsnames]{xcolor}
\definecolor{SectionBlue}{HTML}{1B3A5C}
\definecolor{AccentGray}{HTML}{4A4A4A}
\definecolor{QuoteBorder}{HTML}{8B9DAF}
\definecolor{RiskRed}{HTML}{C0392B}
\definecolor{RiskAmber}{HTML}{E67E22}
\definecolor{RiskGreen}{HTML}{27AE60}
\definecolor{BoxBg}{HTML}{F7F8FA}

% ── Tables ────────────────────────────────────────────────────────────────────
\usepackage{booktabs}
\usepackage{tabularx}

% ── Boxes ─────────────────────────────────────────────────────────────────────
\usepackage[framemethod=tikz]{mdframed}

% ── Lists ─────────────────────────────────────────────────────────────────────
\usepackage{enumitem}
\usepackage{needspace}

% ── Hyperlinks ────────────────────────────────────────────────────────────────
\usepackage{hyperref}
\hypersetup{
  colorlinks=true,
  linkcolor=SectionBlue,
  urlcolor=SectionBlue,
  citecolor=SectionBlue,
  pdfauthor={Pat Singh},
  pdftitle={PR/FAQ — Beadle},
  pdfsubject={Working Backwards — Product Discovery},
}

% ── Bibliography ─────────────────────────────────────────────────────────────
\usepackage[style=numeric,backend=biber,sorting=none]{biblatex}
\addbibresource{prfaq.bib}

% ── Section styling ───────────────────────────────────────────────────────────
\usepackage{titlesec}
\titleformat{\section}
  {\needspace{6\baselineskip}\Large\bfseries\color{SectionBlue}}
  {}
  {0pt}
  {}

\titleformat{\subsection}
  {\needspace{5\baselineskip}\large\bfseries\color{AccentGray}}
  {}
  {0pt}
  {}

\titleformat{\subsubsection}
  {\normalsize\bfseries\color{AccentGray}}
  {}
  {0pt}
  {}
\titlespacing*{\subsubsection}{0pt}{1em}{0.5em}[0pt]

% ── Document stage ───────────────────────────────────────────────────────────
\newcommand{\prfaqstage}[1]{\def\theprfaqstage{#1}}
\prfaqstage{hypothesis}

% ── Document version ─────────────────────────────────────────────────────────
\newcommand{\prfaqversion}[2]{\def\theprfaqmajor{#1}\def\theprfaqminor{#2}}
\prfaqversion{3}{1}

% ── Header/footer ────────────────────────────────────────────────────────────
\usepackage{fancyhdr}
\pagestyle{fancy}
\fancyhf{}
\fancyhead[L]{\footnotesize\textcolor{AccentGray}{Stage: \theprfaqstage{} \enspace\textbar\enspace v\theprfaqmajor.\theprfaqminor}}
\fancyhead[R]{\footnotesize\textcolor{AccentGray}{\thepage}}
\renewcommand{\headrulewidth}{0pt}

% ══════════════════════════════════════════════════════════════════════════════
% Custom environments
% ══════════════════════════════════════════════════════════════════════════════

% ── \prsection{title} — styled press release section header ──────────────────
\newcommand{\prsection}[1]{%
  \vspace{1em}%
  {\large\bfseries\color{SectionBlue}#1}%
  \vspace{0.3em}\par%
}

% ── customerquote — indented, italic customer quote ──────────────────────────
\newmdenv[
  topline=false,
  bottomline=false,
  rightline=false,
  linewidth=3pt,
  linecolor=QuoteBorder,
  backgroundcolor=BoxBg,
  innerleftmargin=12pt,
  innerrightmargin=12pt,
  innertopmargin=8pt,
  innerbottommargin=8pt,
  skipabove=\baselineskip,
  skipbelow=\baselineskip,
]{customerquote}

% ── spokespersonquote — indented spokesperson quote ──────────────────────────
\newmdenv[
  topline=false,
  bottomline=false,
  rightline=false,
  linewidth=3pt,
  linecolor=SectionBlue,
  backgroundcolor=BoxBg,
  innerleftmargin=12pt,
  innerrightmargin=12pt,
  innertopmargin=8pt,
  innerbottommargin=8pt,
  skipabove=\baselineskip,
  skipbelow=\baselineskip,
]{spokespersonquote}

% ── faqpair — numbered Q&A pair with bold question ─────────────────────────
\newcounter{faqnum}
\newenvironment{faqpair}[1]{%
  \needspace{4\baselineskip}%
  \vspace{0.5em}%
  \refstepcounter{faqnum}%
  \noindent\textbf{\color{SectionBlue}Q\thefaqnum: #1}\par\nopagebreak[4]%
  \vspace{0.2em}%
  \begin{adjustwidth}{1em}{0pt}%
  \setlength{\parindent}{0pt}%
}{%
  \end{adjustwidth}%
  \vspace{0.5em}%
}
% ── \faqref{faq:slug} — clickable "FAQ 3" reference ──────────────────────────
\newcommand{\faqref}[1]{\hyperref[#1]{FAQ~\ref*{#1}}}

% ── \featureitem — numbered, referenceable feature entry ──────────────────────
\newcounter{featurenum}
\newcommand{\featureitem}[2]{%
  \refstepcounter{featurenum}%
  \item[\textbf{\color{SectionBlue}F\thefeaturenum.}]%
  \textbf{#1} --- #2%
}
\newcommand{\featureref}[1]{\hyperref[#1]{Feature~\ref*{#1}}}

% ── adjustwidth (for faqpair indentation) ────────────────────────────────────
\usepackage{changepage}

% ══════════════════════════════════════════════════════════════════════════════
% Document
% ══════════════════════════════════════════════════════════════════════════════

\begin{document}

% ── Headline block ───────────────────────────────────────────────────────────
\begin{center}
  {\LARGE\bfseries\color{SectionBlue}
    Punt Labs Launches Beadle} \\[0.8em]
  {\large\color{AccentGray}
    Signed instructions. Declared permissions. Tamperproof audit trail.\\
    Your autonomous agent, on your machine, under your control.}
\end{center}

\vspace{0.5em}

\noindent\textbf{\color{RiskAmber}What's real today, as of v0.16.5:} Everything
from here through the Call to Action --- the whole Press Release
section --- describes Beadle's target end state --- a signed,
autonomous command-pipeline daemon --- written in the Working Backwards
convention, as if that end state already shipped. It hasn't. What ships
today, and is what a user installing Beadle actually gets, is
\texttt{beadle-email}: an MCP email server with a four-level PGP trust
model and per-contact permissions, installed via a single script (see
\href{https://github.com/punt-labs/beadle\#readme}{README.md}). The
pipeline daemon (\texttt{beadle-daemon}) exists as tested, unreleased
code. GPG command-signing enforcement --- the mechanism this document's
headline is named for --- is implemented and wired into the daemon's
command loader, and a fresh, unconfigured deployment runs zero
daemon-triggered commands rather than unsigned ones; what does not yet
exist is a \texttt{beadle init}/\texttt{sign} onboarding flow for an
operator to configure it. The Feature Appendix states the exact split.

\vspace{0.5em}

% ══════════════════════════════════════════════════════════════════════════════
% PRESS RELEASE
% ══════════════════════════════════════════════════════════════════════════════

\section*{Press Release}

SEATTLE, WA --- [target date] --- Punt Labs today released Beadle,
a free, open-source daemon that turns email into a programmable control
plane for an autonomous agent running on your own machine. You send an
email; Beadle verifies your identity via PGP or Proton E2E encryption,
decomposes your instruction into a pipeline of commands, and executes
them --- mixing AI reasoning (Claude) with fast CLI tools (biff, jq,
git) in a single pipeline. Results arrive as an email reply. Unlike
agent frameworks that ask for broad trust, Beadle requires GPG-signed
command definitions, enforces per-contact permissions before executing
anything, and records every step in a pipeline audit trail. Commands
are programs, the daemon is the shell, pipelines are pipes, and GPG
signatures are sudo (see \faqref{faq:competitors}).

\prsection{Problem}

Developers today use AI coding assistants interactively: they sit at the
keyboard, type a prompt, review the output, and iterate. When the laptop closes,
the agent stops. A developer who triggers a two-hour test suite at 6pm must
either keep the terminal open until 8pm or accept that a crash at 7pm means
starting over in the morning. Delegating multi-step tasks --- run tests across
three environments, summarize the results, file issues for failures --- means
cobbling together fragile cron jobs and shell scripts with no authentication, no
permissions model, and no audit trail. Existing autonomous agent frameworks like Auto-GPT run locally but offer
no cryptographic verification of who authorized which action, no per-command
permissions, and no tamperproof record of what happened.
The risk is not theoretical: in January 2026, ClawdBot (OpenClaw) went
viral, growing from roughly 9,000 to over 60,000 GitHub stars in 72 hours
and surpassing 145,000 within about two weeks \cite{cnbc2026openclawstars},
and was immediately found to have critical vulnerabilities --- authentication
disabled by default, a one-click remote-code-execution chain, and a
spoofed-email attack that exfiltrated API keys
\cite{kaspersky2026_openclaw, socradar2026cve25253, snyk2026_clawdbot}.
Autonomous agents without enforceable trust boundaries produce real security
failures at scale (see \faqref{faq:why-now}).

\prsection{Solution}

Beadle is a background daemon with its own email address, GPG keychain,
and a library of GPG-signed YAML command files. Each command declares
a runner (Claude for AI reasoning, CLI for deterministic tools), its
inputs and outputs via JSON schema, and whether it transforms the pipeline
data or passes it through as a side-effect. Commands compose like Unix
pipes:

{\small\texttt{summarize | notify-team | reply}}

The owner sends an email to Beadle's address. The daemon verifies the
sender's identity (PGP signature or Proton E2E encryption), checks their
\texttt{x} (execute) permission, maps the instruction to a command
pipeline, and executes stages sequentially. A single JSON value --- the
pipe --- flows through the pipeline. Process-mode commands transform it;
passthrough-mode commands (notifications, logging) read it without
modification. Results arrive as an email reply. Every stage is logged
with pipeline ID, command name, runner, duration, and exit status
(see \faqref{faq:why-email}).

\prsection{Customer Quote}

\begin{customerquote}
\textit{``I emailed `summarize this' from my phone on the train and had a
structured summary back in my inbox in under two minutes --- reasoning
done, the headline already broadcast to my team, before I sat down at my
desk. I didn't touch a laptop. When my manager later asked whether the
agent had actually sent that update or I'd done it myself, I could point
at the audit log and answer with certainty instead of a guess.''}

\raggedleft --- \textbf{Alex Petrov}, Senior Platform Engineer, mid-size
fintech
\end{customerquote}

\prsection{Getting Started}

Getting started takes three steps. First, install the binary:
{\small\texttt{curl -fsSL .../install.sh | sh}} downloads the Go binary
and installs the Claude Code plugin (MCP tools + slash commands). Second,
configure your email connection (IMAP/SMTP via Proton Bridge or any
provider) and store credentials in your OS keychain. Third, start the
daemon: {\small\texttt{beadle-daemon run}}. Beadle polls your inbox,
and when you send an email with ``summarize'' in the subject, you get
a structured summary back within two minutes. No cloud account, no data
leaves your machine. Commands are YAML files in
{\small\texttt{\textasciitilde/.punt-labs/beadle/commands/}} --- add new
capabilities by writing a YAML file and signing it with your GPG key.

\prsection{Spokesperson Quote}

\begin{spokespersonquote}
``At Punt Labs we build on a principle: earn trust to go fast.
Every autonomous agent framework asks you to trust it implicitly --- give it
your credentials, let it run whatever it wants, and hope the logs are honest.
We built Beadle on the opposite assumption: the agent has zero authority
of its own. Every signed instruction, every declared permission, every
tamperproof log entry earns the right to automate more aggressively tomorrow.
The signing ceremony isn't overhead --- it's the mechanism that compounds trust
so you can hand off bigger tasks with confidence. It's the Unix philosophy
applied to AI agents: small, composable, auditable tools that do one thing
well.''

\raggedleft --- \textbf{Pat Singh}, Punt Labs
\end{spokespersonquote}

\prsection{Call to Action}

Beadle is available now at
{\small\texttt{github.com/punt-labs/beadle}} under the MIT license.
Static Go binaries for macOS and Linux (arm64/amd64). Installation
requires an IMAP/SMTP server (Proton Bridge recommended for E2E
encryption) and GPG for command signing. Documentation, example
commands, and the pipeline design document are in the repository.

\newpage

% ══════════════════════════════════════════════════════════════════════════════
% FREQUENTLY ASKED QUESTIONS
% ══════════════════════════════════════════════════════════════════════════════

\section*{Frequently Asked Questions}

% ── External FAQs (Customer-facing) ──────────────────────────────────────────

\subsection*{External FAQs}

\begin{faqpair}{What is Beadle and who is it for?}\label{faq:what-is-it}
  Beadle is an open-source agent daemon that receives instructions via email
  and executes them as multi-stage pipelines. It is for developers who want
  to delegate tasks to an AI agent running on their own machine --- with
  cryptographic verification of who authorized what. Commands are YAML files
  that declare a runner (Claude for reasoning, CLI for tools like jq, biff,
  git), a JSON output schema, and whether they transform the pipeline data
  or pass it through as a side-effect. The daemon is the shell; commands
  are programs; GPG signatures are sudo.
\end{faqpair}

\begin{faqpair}{How is Beadle different from Auto-GPT or Claude Code?}\label{faq:competitors}
  Auto-GPT and similar frameworks give an agent broad autonomy with minimal
  authentication. Claude Code is an interactive assistant that requires the
  developer at the keyboard (its \texttt{-{}-bare} mode enables headless
  execution but provides no trust model, no command library, and no pipeline
  orchestration).

  Beadle occupies a different point: it has \textit{zero} independent authority.
  Every command is a GPG-signed YAML file with a declared runner, typed args,
  and JSON output schema. Pipelines compose explicitly. The trust gate verifies
  sender identity (PGP or Proton E2E) and per-contact permissions before any
  execution. CLI commands run in milliseconds with whitelist-only binary
  resolution and no shell. The trade-off is deliberate: Beadle sacrifices agent
  autonomy for owner control, auditability, and speed on deterministic operations.
\end{faqpair}

\begin{faqpair}{Why email? Why not a web UI or a CLI?}\label{faq:why-email}
  Email is asynchronous, universal, and already authenticated. Proton Mail
  provides end-to-end encryption. GPG provides sender verification. Together,
  they give Beadle a control plane that works from any device --- phone, tablet,
  another computer --- without installing a client or opening a port. The
  concept of email as an agent control channel has engineering precedent:
  HumanLayer's Agent Control Plane uses email for human-approval of autonomous
  agent actions \cite{humanlayer2024acp}. The owner
  can CC Beadle on a conversation with a colleague, and Beadle can distinguish
  the owner's signed instructions from the colleague's unsigned replies.

  The vision includes local execution via a CLI
  ({\small\texttt{beadle run pipeline.yaml}}) for development and testing.
  Email is meant to be the remote interface, not the only one.
\end{faqpair}

\begin{faqpair}{How do I get started?}\label{faq:getting-started}
  This is the target flow, not what's installable today (see
  \faqref{faq:timeline} for the gap): run the install script
  ({\small\texttt{curl -fsSL .../install.sh | sh}}) to download the Go
  binary and install the Claude Code plugin. Configure your email
  connection (IMAP/SMTP) and credentials. Start the daemon with
  {\small\texttt{beadle-daemon run}}. Send an email with ``summarize'' in
  the subject from a PGP-verified or Proton E2E contact, and receive the
  summary as a reply. Add commands by writing YAML files in
  {\small\texttt{\textasciitilde/.punt-labs/beadle/commands/}}.
\end{faqpair}

\begin{faqpair}{How much does it cost?}\label{faq:pricing}
  Beadle is free and open-source. The recommended path is a Proton Mail
  paid plan (starting at \$3.99/month), because Proton Bridge --- the
  local IMAP/SMTP server this setup is built around --- is not available
  on the free tier \cite{protonbridge2025}. It is not the only path: Beadle
  connects over standard IMAP/SMTP with no Proton-specific dependency (see
  \faqref{faq:tech-risks}), and Fastmail is already used for the project's
  own signing tests. You also need a GPG keychain (free). If your command
  pipelines use the Claude interpreter, you will need an Anthropic API key
  with usage-based billing --- unquantified in this document; see
  \faqref{faq:pnl} for the corresponding gap on the Punt Labs side. Beadle
  itself has no cost.
\end{faqpair}

\begin{faqpair}{What happens to my data?}\label{faq:data}
  Nothing leaves your machine unless a command document explicitly sends it
  somewhere. Beadle runs locally as a daemon. With the recommended Proton
  Bridge setup, email flows through Bridge on your machine to Proton's
  encrypted servers; with any other IMAP/SMTP provider, it flows to that
  provider's servers instead (see \faqref{faq:pricing}). Command documents,
  audit logs, and results stay on your local filesystem. There is no
  telemetry, no cloud service, and no account with Punt Labs.
\end{faqpair}

\begin{faqpair}{Do I need to know GPG?}\label{faq:gpg-knowledge}
  Today: not for identity. Beadle resolves who you are through
  \texttt{ethos}, an identity sidecar shared across Punt Labs tools --- and
  if you don't want even that, one plain-text file
  ({\small\texttt{\textasciitilde/.punt-labs/beadle/default-identity}}) is
  enough. Signing uses your own GPG key, configured once during
  {\small\texttt{beadle-email install}}; you never invoke \texttt{gpg}
  directly, and the key must carry an expiration date --- Beadle already
  rejects non-expiring keys, today, for this signing.

  For command signing specifically --- the mechanism this document's
  headline is named for --- the honest answer is different: it isn't
  built, so there is no onboarding flow to evaluate yet. The dedicated
  {\small\texttt{beadle init}}/{\small\texttt{sign}}/{\small\texttt{verify}}
  commands imagined in earlier drafts of this document were never built;
  what shipped instead solves identity, not command-signing enforcement.
  When enforcement ships, it needs its own onboarding design --- a Must
  Do, not yet started (see \faqref{faq:timeline}).
\end{faqpair}

% ── Internal FAQs (Business-facing) ──────────────────────────────────────────

\newpage
\subsection*{Internal FAQs}

\subsubsection*{Value \& Market}

\begin{faqpair}{What is the total addressable market?}\label{faq:tam}
  The primary market is developers who already use AI coding assistants and want
  unattended execution. LangChain's November--December 2025 survey (1,340
  practitioners) found 57.3\% have agents in production, up from 51\% the
  year before \cite{langchain2025agenteng, langchain2024agents}.
  The subset who need local-first, authenticated, autonomous execution is
  narrower --- likely tens of thousands of developers who work with sensitive
  codebases, regulated environments, or privacy-conscious infrastructure
  (the Proton Mail user base is an imperfect proxy for this audience).

  This is not a rigorous bottoms-up estimate --- there's no segment count
  or willingness-to-pay analysis behind it, just a top-down survey stat and
  an imperfect proxy. We are not claiming a billion-dollar TAM. The honest
  framing: Beadle targets a niche of security-conscious developers who
  value auditability over convenience, sized by inference, not
  arithmetic.
\end{faqpair}

\begin{faqpair}{What evidence do we have that customers want this?}\label{faq:customer-evidence}
  At hypothesis stage, the evidence is almost entirely internal. The
  full pipeline runs reliably against the author's own mailbox --- a
  real, PGP-verified email triggers a multi-stage pipeline (Claude and
  CLI runners, JSON schema validation between stages) and returns a
  result by reply. That validates the mechanics. It does not validate
  that anyone besides the person who built it wants it: zero external
  developers have used Beadle (see \faqref{faq:next-step}).

  \textbf{Inferential evidence:}
  \begin{itemize}[nosep]
    \item Auto-GPT has accumulated over 187,000 GitHub stars since its 2023
      release, demonstrating strong developer interest in autonomous agents
      \cite{autogpt2024github}.
    \item Cisco's 2025 Data Privacy Benchmark Study found 64\% worry about
      inadvertently sharing sensitive information via GenAI tools even as
      90\% see local storage as inherently safer --- though the same study
      found 91\% trust global providers for better data protection than
      they could provide themselves, up 5 points year-over-year. The market
      is genuinely split, not uniformly local-preferring
      \cite{cisco2025privacy}.
    \item The OWASP GenAI Security Project published a Top~10 for Agentic
      Applications in 2025, indicating that agent security is a recognized
      concern, not a fringe worry \cite{owasp2025agentic}.
    \item Proton Mail has grown to 100M+ accounts, confirming demand for
      privacy-focused communication infrastructure \cite{proton2024100m}.
    \item No existing agent framework offers GPG-based authentication with
      per-command permissions --- the gap is observable, but unvalidated.
  \end{itemize}

  \textbf{Next validation step:} see \faqref{faq:next-step} --- it is no
  longer ``generalize the prototype.'' The prototype has been generalized.
  What's missing is distribution and the onboarding flow that would let
  an external developer reach it at all.
\end{faqpair}

\begin{faqpair}{Who are the competitors and why will we win?}\label{faq:competitive-landscape}
  \begin{tabularx}{\linewidth}{@{} l X X @{}}
  \toprule
  \textbf{Competitor} & \textbf{Strength} & \textbf{Trust model gap} \\
  \midrule
  Auto-GPT & Large community, broad agent autonomy &
    Agent-directed execution; no owner authentication, no declared
    permissions, no tamperproof audit trail \\
  Claude Code & High-quality code generation, interactive &
    Operator-attended trust model; requires user at keyboard,
    cannot delegate unattended authority \\
  Devin & Full autonomous dev environment &
    Vendor-hosted trust model; execution runs on Cognition's
    infrastructure, not the owner's machine \\
  Cron + scripts & Simple, well-understood &
    No trust model; no identity verification, no permission
    scoping, no composability with LLMs \\
  \bottomrule
  \end{tabularx}

  \vspace{0.5em}
  Beadle does not compete on agent intelligence --- it uses existing LLMs
  (Claude, GPT) as interpreters. The differentiation is the trust model:
  local-only execution under cryptographic owner control. Cloud providers and
  SaaS agent platforms cannot replicate this without contradicting their own
  business model --- their revenue depends on hosting execution, managing
  infrastructure, and retaining custody of the runtime environment. Offering
  genuine local-only, owner-controlled execution means giving up the recurring
  billing, telemetry, and platform lock-in that sustain their economics.
  Competitors can add daemon modes, unattended execution, and even signing
  features, but they cannot cede control of the execution environment to the
  owner without dismantling the value capture that funds their products. This
  is a structural misalignment, not a feature gap.
\end{faqpair}

\begin{faqpair}{What is the customer acquisition strategy?}
  Open-source distribution via GitHub Releases. Beadle is a CLI tool in the Punt
  Labs ecosystem (alongside Biff, Quarry, and Vox), which provides cross-project
  visibility. The secondary purpose --- as a showcase for the Jacobson-Cockburn
  use-case methodology \cite{jacobson2023usecases} --- generates content that
  attracts developers interested in both the tool and the process.

  No paid acquisition. Growth is organic: one launch post (Hacker News,
  Lobsters), 2--3 blog posts documenting the build process and methodology,
  and cross-links from other Punt Labs project READMEs. Estimated effort:
  20--30 hours of content creation in the first 3 months. If that generates
  fewer than 100 stars, the positioning missed the audience and content
  investment should stop.
\end{faqpair}

\begin{faqpair}{What is your next step to validate your vision?}\label{faq:next-step}
  Both Pipeline v1 and Pipeline v2 shipped as tested code since this
  document's v3.0 draft --- CLI runners, process/passthrough data flow,
  compound CLI steps, and JSON schema validation all merged in
  \texttt{v0.15.0} (2026-04-18). \texttt{beadle-email} itself has shipped
  six releases since, through \texttt{v0.16.5}. Building has not been
  the bottleneck.

  Validation has. The external test proposed at v3.0 --- 5--10 developers
  who use Claude Code or Auto-GPT for automation, writing a command YAML,
  signing it, sending an email, and getting a pipeline result back --- has
  not run, and cannot run yet: \texttt{beadle-daemon} is not distributed to
  anyone (\faqref{faq:tech-risks} covers why), GPG command-signing
  enforcement is an unimplemented stub, and there is no
  {\small\texttt{beadle init}} to abstract key setup.

  The next step is therefore not more pipeline features. In order: (1) ship
  GPG signing enforcement, so the trust model this document is named for is
  real; (2) build the {\small\texttt{beadle init}} onboarding flow; (3)
  distribute \texttt{beadle-daemon}; then (4) finally run the external
  test. Until that test runs, the riskiest assumption in this document ---
  whether the signing ceremony feels like a feature or a burden --- is
  exactly as untested as it was at v3.0.
\end{faqpair}

\subsubsection*{Technical}

\begin{faqpair}{What are the major technical risks?}\label{faq:tech-risks}
  \begin{enumerate}[nosep]
    \item \textbf{Proton Bridge reliability for automated use.} Proton Bridge
      is designed for desktop email clients, not daemons. It may not handle
      long-running, unattended operation gracefully.

      \textit{Why Proton specifically:} Beadle's threat model assumes the
      email channel must protect command integrity end-to-end --- from the
      owner's device to the daemon's mailbox. Proton Mail's end-to-end
      encryption means even Proton's own servers cannot read command content in
      transit, which aligns with the local-first, zero-trust design. That said,
      Beadle connects to Proton Bridge via standard IMAP/SMTP --- there is no
      Proton-specific API or protocol dependency. Supporting other email
      providers (Fastmail, self-hosted Dovecot, Gmail with app passwords) is a
      scope decision, not an architectural constraint.

      \textit{Design, not yet built:} Beadle should monitor Bridge health on
      a configurable interval and expose status via a status command; on
      outage it should log, queue pending result deliveries, and retry on a
      backoff schedule, with a service-manager restart (macOS
      {\small\texttt{LaunchAgent}} + {\small\texttt{KeepAlive}}, or a systemd
      unit with {\small\texttt{Restart=on-failure}}) if the daemon itself
      exits. None of this exists today. \texttt{beadle-daemon} has no
      health monitor, retry queue, status command, or service-manager
      integration --- if Bridge goes down while the daemon is running,
      email-triggered execution simply stops with no recovery mechanism.
    \item \textbf{GPG key management complexity.} GPG tooling is notoriously
      difficult --- complexity is consistently cited as its primary operational
      weakness \cite{hoopdotdev2024gpg}. If {\small\texttt{beadle init}} cannot
      fully abstract key generation and signing, the onboarding friction will
      kill adoption. This is the hardest unsolved UX problem in the project.

      \textit{Alternative evaluated: SSH signing.} GitHub and GitLab both
      support SSH keys for commit signing, and most developers already have an
      SSH key. However, SSH signing only covers commit-level attestation ---
      it provides no standard for encrypting documents to a recipient, no
      built-in key expiry or revocation model, and no ecosystem for
      distributing trust (no equivalent of the OpenPGP Web of Trust or
      keyservers). Beadle needs signing \textit{and} encryption \textit{and}
      key lifecycle management; SSH covers only the first.

      \textit{What's actually built:} identity resolution already works,
      though not through a bespoke {\small\texttt{beadle init}} flow --- it
      goes through \texttt{ethos}, an identity sidecar shared across Punt
      Labs tools. A user with no ethos installed at all still gets a working
      identity from one plain-text file
      ({\small\texttt{\textasciitilde/.punt-labs/beadle/default-identity}}).
      Signing uses the user's own GPG key from the system keyring ---
      not a Beadle-generated one, and not an isolated keychain; only
      inbound signature \textit{verification} runs in an isolated
      {\small\texttt{GNUPGHOME}}, since that's where an untrusted
      sender's key needs to be quarantined. The signer is configured via
      a {\small\texttt{gpg\_signer}} field the install flow already
      prompts for, and Beadle rejects non-expiring keys before every
      sign. The user never types a \texttt{gpg} command --- but does need
      to already own an expiring key and place its passphrase in the
      keychain by hand during {\small\texttt{beadle-email install}}; that
      residual friction is exactly what {\small\texttt{beadle init}}
      exists to remove.

      \textit{What's still vision:} dedicated {\small\texttt{beadle
      init}}/{\small\texttt{sign}}/{\small\texttt{verify}} subcommands. The
      daemon's command-signature verifier ({\small\texttt{VerifySignature}},
      DES-034) is implemented --- canonical signing, isolated verification
      against a known command-file authorizer key, GnuPG status-fd classification
      --- and wired into the command-file loader (DES-035), so the daemon
      rejects an unsigned or tampered command file, and starts with zero
      loadable commands until an operator sets {\small\texttt{owner\_handle}}
      or {\small\texttt{owner\_gpg\_key\_id}} in {\small\texttt{daemon.json}}
      (a fresh, unconfigured deployment runs no daemon-triggered commands at
      all, never unsigned ones). There is still no {\small\texttt{beadle
      init}} onboarding to make that configuration step easy, or a dedicated
      signing tool for a new user to sign their own command files with ---
      that onboarding gap is now the higher-priority remaining piece (see
      \faqref{faq:next-step}).
    \item \textbf{Claude interpreter sandboxing.} When Beadle invokes Claude
      as a command interpreter, the LLM may attempt actions outside the
      command's declared permissions --- or a message's \textit{content},
      not just its sender, could carry a prompt injection, exactly the
      vector Snyk demonstrated against ClawdBot (see \faqref{faq:why-now}).
      \textit{Mitigation, already built:} workers spawn via
      {\small\texttt{claude -p --bare}}, which skips CLAUDE.md, hooks,
      plugins, and auto-memory; an adversarial system prompt frames every
      worker invocation; the worker's \texttt{PATH} is restricted to the
      Claude binary directory plus \texttt{/usr/bin}; and the daemon
      withholds its own credentials from the worker's environment.
      \textit{Still open:} permission preflighting is authorization
      bookkeeping, not a content firewall --- it establishes who sent a
      command, not what a message's body asks the interpreter to do once
      inside its declared scope, and the worker still runs with the real
      {\small\texttt{\$HOME}} and broad tool access
      ({\small\texttt{Bash}}, {\small\texttt{Edit}}, {\small\texttt{Write}}).
      There is no container-level isolation (network namespace, seccomp)
      yet.
  \end{enumerate}
\end{faqpair}

\begin{faqpair}{What dependencies exist on other teams or systems?}
  \begin{itemize}[nosep]
    \item \textbf{Proton Bridge} --- open-source Go codebase, maintained by
      Proton AG \cite{protonbridge2025}. Beadle uses it as a standard
      IMAP/SMTP client. No API dependency beyond the mail protocol.
    \item \textbf{Anthropic API} --- required only when using the Claude
      interpreter. Usage-based billing, no contractual dependency.
    \item \textbf{GnuPG} --- the GPG implementation. Mature, widely deployed,
      maintained by the GnuPG Project. RFC~9580 (OpenPGP, 2024) provides
      the current standard \cite{rfc9580}.
  \end{itemize}
  No internal Punt Labs dependencies. Beadle is a standalone project.
\end{faqpair}

\begin{faqpair}{What is the estimated development timeline?}\label{faq:timeline}
  This is a hypothesis-stage document. The purpose of Working Backwards is
  to validate \textit{whether} to build, not to plan the schedule. Detailed
  hour estimates and calendar targets are premature --- and doubly so when
  AI-augmented development makes historical velocity assumptions unreliable.
  What matters at this stage is delivery order and cut discipline.

  \textbf{Phased delivery order:}

  \begin{enumerate}[nosep]
    \item \textbf{Foundation (shipped, distributed)} --- MCP email server
      with 21 always-on tools plus 2 poll tools gated on config,
      four-level PGP trust model, per-contact rwx
      permissions, multi-identity via ethos (with a plain-file fallback
      that needs no ethos at all), credential isolation (OS keychain).
      Released as \texttt{beadle-email}, currently \texttt{v0.16.5}. This
      is the only layer any user can actually install today.
    \item \textbf{Daemon + Pipeline v1 (built, not distributed)} ---
      Background daemon with IMAP poller, trust gate (PGP verified +
      Proton E2E), x-bit enforcement, pipeline orchestrator (command YAML
      loader, regex planner, sequential executor, auto-reply, crash-safe
      persistence), worker spawner via \texttt{claude -p --bare} with
      prompt-level adversarial hardening (see \faqref{faq:tech-risks} for
      what that does and doesn't cover). Tested and e2e verified against a
      real mailbox, but \texttt{make dist} and \texttt{install.sh} do not
      build or ship it.
    \item \textbf{Pipeline v2 (built, not distributed)} --- CLI runners
      (execute tools in milliseconds instead of a full Claude session),
      the process/passthrough data-flow model, compound CLI steps
      (jq | biff without a shell), JSON pipe with schema validation,
      runner-conditional command validation. Shipped as tested code in
      \texttt{v0.15.0} (2026-04-18) --- the same distribution gap as
      Pipeline v1 applies.
    \item \textbf{Future (not built)} --- GPG command-signing enforcement
      ({\small\texttt{VerifySignature}} is implemented, DES-034, but not
      wired into the command-file loader --- see \faqref{faq:next-step}),
      the {\small\texttt{beadle init}} onboarding flow, daemon distribution,
      an LLM planner (replace regex with Claude decomposition),
      multi-channel triggers (biff chat, webhooks), an internal scheduler,
      key rotation management.
  \end{enumerate}

  \textbf{Cut order:} Future items are cut first, but the first three ---
  signing enforcement, {\small\texttt{beadle init}}, and distribution ---
  are what stand between this document's own thesis and a product a user
  can actually evaluate. Everything after those three is genuinely
  deferrable; those three are not. The kill switch is adoption-based: see
  \faqref{faq:metrics}.
\end{faqpair}

\begin{faqpair}{What is the scaling story?}\label{faq:scaling}
  Beadle is a single-user daemon. ``Scaling'' means long-running operation over
  months and years, not multi-tenant growth.

  \begin{itemize}[nosep]
    \item \textbf{Audit log growth.} The log is append-only and unbounded until
      the owner manually clears it. A developer running 10 pipelines per day
      generates roughly 50--100~KB of log per day. At one year, that is
      approximately 20--40~MB --- manageable on any modern filesystem. The risk
      is neglect, not size.
    \item \textbf{Command library size.} Commands are loaded into memory only
      during preflight and execution, then discarded. A library of hundreds of
      signed command documents on disk has no runtime cost until invoked.
    \item \textbf{Proton Bridge longevity.} The daemon depends on Bridge staying
      online for email-triggered execution. Bridge was not designed for months
      of continuous uptime. \textit{Design, not yet built:} the mitigation
      should be operational --- Beadle should check Bridge health on a
      configurable interval, a service manager should restart both Beadle
      and Bridge on failure, and pending result deliveries should queue and
      retry once Bridge recovers. None of this exists today (see
      \faqref{faq:tech-risks}, \featureref{feat:bridge-health}); an
      internal scheduler that would let non-email-triggered pipelines run
      independently of Bridge is also unbuilt (\featureref{feat:scheduler}).
  \end{itemize}

  The honest constraint: Beadle does not scale to teams, organizations, or
  multi-machine deployment. It is a personal agent for one person on one
  machine. That is a design choice, not a limitation to fix later.
\end{faqpair}

\subsubsection*{Business}

\begin{faqpair}{What is the revenue model?}\label{faq:revenue}
  None. Beadle is free and open-source. The strategic value is threefold:
  (1)~it demonstrates the Punt Labs engineering methodology in a real,
  useful product; (2)~it grows the Punt Labs open-source ecosystem alongside
  Biff, Quarry, and TTS; (3)~it serves as the reference implementation for
  a future punt-labs project on use-case methodology, generating content and
  credibility.
\end{faqpair}

\begin{faqpair}{What does the P\&L look like at steady state?}\label{faq:pnl}
  Beadle has no revenue, so the P\&L is cost-only. The primary cost is
  developer-hours. Specific \textit{development}-hour estimates are not
  attempted at hypothesis stage --- the purpose of this document is to
  validate whether to build,
  not to plan the schedule --- but the cost categories are known.

  \begin{tabularx}{\linewidth}{@{} l X @{}}
  \toprule
  \textbf{Cost} & \textbf{Notes} \\
  \midrule
  Development (primary) & Solo developer. Four phases described in
    \faqref{faq:timeline}. The dominant cost and the binding constraint
    on opportunity cost. \\
  Ongoing maintenance & Bug fixes, dependency updates, community PRs.
    Modest but nonzero indefinitely. \\
  Proton Mail (paid plan) & \$48/year. Personal use; already paid for
    other purposes. \\
  Infrastructure & \$0. No servers; runs on developer's machine. \\
  \bottomrule
  \end{tabularx}

  \vspace{0.5em}
  The opportunity cost is the binding constraint: development time spent on
  Beadle is time not spent on other Punt Labs projects (Quarry features, Biff
  improvements, LangLearn ports). This cost is governed by a kill switch tied
  to the metrics in \faqref{faq:metrics}: if, at the 6-month mark after
  public release, Beadle has fewer than 100 GitHub stars \textit{or} fewer
  than 10 weekly binary downloads, the project enters archival mode (security
  patches only, no new feature work). This is a hard threshold, not a
  judgment call --- below either number means the positioning missed the
  audience and further development hours are not justified.

  \textbf{Where the clock actually stands:} first public release
  ({\small\texttt{v0.1.0}}) was 2026-03-13. The 6-month mark is
  2026-09-13 --- about two weeks from this revision's date. The current
  star count is 3. That is not ``below the 100-star failure threshold'';
  it is roughly 3\% of it. The kill switch as originally written also
  measures the wrong thing: \texttt{beadle-daemon}, the product this
  document is about, isn't distributed (see \faqref{faq:timeline}), so
  ``weekly binary downloads'' can only mean \texttt{beadle-email}
  downloads --- a proxy for interest in the email MCP server, not in the
  signed-pipeline-daemon vision.
\end{faqpair}

\begin{faqpair}{What are the key metrics and how will we measure success?}\label{faq:metrics}
  \begin{tabularx}{\linewidth}{@{} l l X @{}}
  \toprule
  \textbf{Metric} & \textbf{Target (6 months)} & \textbf{Decision threshold} \\
  \midrule
  GitHub stars & 500+ & Below 100 after launch suggests the positioning
    missed the audience \\
  Weekly \texttt{beadle-email} downloads & 50+ & Below 10 suggests the
    email MCP server, specifically, is not useful enough to keep \\
  External pipeline test completion & 8/10 users complete first pipeline &
    Below 5/10 signals fatal usability friction --- untested; blocked on
    \faqref{faq:next-step} \\
  \bottomrule
  \end{tabularx}

  \vspace{0.5em}
  \textbf{These metrics don't measure the actual product.} Star count and
  \texttt{beadle-email} downloads track interest in the email MCP server
  Punt Labs already ships. They say nothing about the signed
  command-pipeline daemon --- the thing this document's headline, press
  release, and competitive positioning are all about --- because that
  daemon has never been distributed to anyone. A revised kill switch needs
  either its own instrumentation for the daemon once it ships, or an
  honest statement that the current metrics are a proxy for a different,
  adjacent hypothesis (``does an MCP email server with a trust model have
  users?'') than the one this document exists to test.

  \textbf{Recommendation:} do not run the kill switch as originally
  written --- it would fire on a metric that was never measuring the
  right thing. Re-baseline the clock to \texttt{beadle-daemon}'s actual
  distribution date (see \featureref{feat:daemon-dist}), and hold the
  Must Do items to their own accountability instead: if signing
  enforcement, {\small\texttt{beadle init}}, and daemon distribution
  aren't shipped within one quarter of this revision, that --- not a
  star count --- is the signal to archive.

  \textbf{Where the clock actually stands, as of this revision:} 3 GitHub
  stars against a first-release date of 2026-03-13, with the 6-month mark
  landing 2026-09-13 --- about two weeks out. Whatever the right
  kill-switch design turns out to be, the current trajectory does not
  clear even the ``below 100 means it missed'' failure bar, let alone the
  500-star target.

  \textbf{Pre-mortem --- what failure looks like:} written at v3.0,
  imagining March 2027. The failure condition it describes could now
  arrive well before then, for the reason stated above rather than the
  one originally imagined.

  It is March 2027. Beadle has
  200 GitHub stars, 8 binary downloads per week, and the prototype test showed
  3/10 users abandoned during GPG key setup. The cryptographic trust model ---
  signed commands, declared permissions, tamperproof audit trail --- resonated
  in blog posts but not in practice. Developers found the signing ceremony
  too heavy for personal automation tasks where they already trust themselves.
  The security posture Beadle enforces felt like overhead, not protection.
  The project becomes a methodology showcase only, with no active user
  community, and enters archival maintenance.
\end{faqpair}

\begin{faqpair}{Why now? What has changed?}\label{faq:why-now}
  Four shifts make this the right time:

  \begin{enumerate}[nosep]
    \item \textbf{AI agents went mainstream in 2024--2025.} LangChain's 2024
      survey found 51\% production deployment \cite{langchain2024agents}; its
      November--December 2025 follow-up (1,340 practitioners) puts that at
      57.3\%, still climbing \cite{langchain2025agenteng}. Developers are past
      the question of ``should I use an agent?'' and into ``how do I trust
      one?''
    \item \textbf{OpenPGP modernized.} RFC~9580 (2024) updated the OpenPGP
      standard for the first time in decades, adding modern cryptographic
      algorithms and improving key management \cite{rfc9580}. GPG tooling is
      better than it was even two years ago.
    \item \textbf{OWASP formalized agent security risks.} The OWASP Top~10 for
      Agentic Applications (2025) legitimized the concern that autonomous agents
      need security models, not just capability models
      \cite{owasp2025agentic}. Beadle's trust model is a direct response.
    \item \textbf{ClawdBot proved the threat model in production.} In January
      2026, ClawdBot (also known as OpenClaw) went viral, growing from
      roughly 9,000 to over 60,000 GitHub stars in 72 hours and surpassing
      145,000 within about two weeks \cite{cnbc2026openclawstars}. Within
      days, security researchers found it riddled with critical
      vulnerabilities. CVE-2026-25253 (CVSS~8.8) was a one-click RCE: the
      Control UI trusted an unvalidated \texttt{gatewayUrl} query parameter,
      auto-connected a WebSocket to it, and leaked the operator's gateway
      auth token in the connect payload --- enabling cross-site WebSocket
      hijacking, sandbox disable, and Docker escape to full host compromise
      from a single malicious link \cite{socradar2026cve25253,
      thn2026openclawrce}. CVE-2026-24763 was a separate, authenticated
      command-injection flaw via unsafe \texttt{PATH} handling in Docker
      sandbox execution \cite{ghsa2026mc68clawdocker}. Separately again, a
      February 2026 audit found 341 malicious skills among the 2,857
      published on OpenClaw's skill marketplace (11.9\% of the registry),
      most traced to one campaign distributing infostealer malware to an
      estimated 300,000 users \cite{koisecurity2026clawhub}. Snyk
      demonstrated that a spoofed email could trigger a shell command that
      exfiltrated \texttt{clawdbot.json} containing API keys
      \cite{snyk2026_clawdbot}. Palo Alto Networks' Unit~42 concluded that
      OpenClaw ``maps to all 10 OWASP risks for agentic applications'' and
      lacked ``enforceable trust boundaries between untrusted inputs and
      high-privilege reasoning'' \cite{unit42_2026_openclaw}. Kaspersky
      reported that authentication was disabled by default
      \cite{kaspersky2026_openclaw}.

      Of these, only the Snyk finding is squarely the failure mode Beadle's
      trust model addresses: a spoofed sender triggering unauthenticated
      execution is exactly what per-contact permissions and PGP
      sender-verification prevent. The gateway RCE was a browser-mediated
      credential-theft chain and the skill-marketplace compromise was a
      supply-chain vetting failure; Beadle's design doesn't claim to address
      either. The honest framing is narrower than ``exactly the failure mode
      Beadle prevents'': OpenClaw proved that shipping an agent with
      \textit{any} weak trust boundary is dangerous at scale, and Beadle's
      differentiator is the one boundary --- sender authentication --- that
      incident specifically validates.
  \end{enumerate}
\end{faqpair}

% ══════════════════════════════════════════════════════════════════════════════
% FOUR RISKS ASSESSMENT
% ══════════════════════════════════════════════════════════════════════════════

\newpage
\section*{Risk Assessment}

\renewcommand{\arraystretch}{1.6}
\begin{tabularx}{\textwidth}{@{} l l X @{}}
\toprule
\textbf{\color{SectionBlue}Risk} & \textbf{\color{SectionBlue}Rating} & \textbf{\color{SectionBlue}Assessment} \\
\midrule
Value & \textcolor{RiskAmber}{Medium} &
  Developer interest in autonomous agents is proven (Auto-GPT stars,
  LangChain survey). Whether the cryptographic trust model --- signed commands,
  declared permissions, tamperproof audit trail --- resonates with developers
  who already trust their own machines is the core untested assumption.
  The security posture may feel like overhead for personal automation.
  Validation plan: prototype test with 5--10 developers
  (see \faqref{faq:next-step}). \\
Usability & \textcolor{RiskRed}{High} &
  GPG abstraction is the hardest unsolved UX problem in the project. The
  {\small\texttt{beadle init}} command must generate keys and produce an
  identity file without ever exposing raw GPG semantics to the user. SSH
  signing was evaluated and rejected because it lacks encryption and key
  lifecycle primitives (see \faqref{faq:tech-risks}). If the abstraction
  leaks --- if any step requires the user to understand keyrings, trust
  models, or \texttt{gpg} flags --- onboarding friction will kill the
  product. The signing ceremony (sign every command, declare every
  permission) is the trust model's core value but also its primary
  usability risk: developers may perceive it as ceremony, not safety. \\
Feasibility & \textcolor{RiskAmber}{Medium} &
  Pipeline v1 and v2 are both built and e2e verified: email arrival
  triggers a multi-stage pipeline (Claude and CLI runners, JSON schema
  validation) and delivers results to the sender. The trust gate, worker
  spawner, pipeline executor, and crash-safe persistence all work against
  a real mailbox. What's untested is the opposite of ``more code'':
  Proton Bridge has run continuously for weeks, not the months this
  document's own scaling FAQ (\faqref{faq:scaling}) says is the real bar,
  and the primary Usability mitigation ({\small\texttt{beadle init}}) is
  entirely unwritten. The hardest remaining work is distribution and
  onboarding, not pipeline mechanics. \\
Viability & \textcolor{RiskAmber}{Medium} &
  No revenue requirement, no infrastructure cost beyond the developer's
  own machine --- on that axis the risk is genuinely low. But the stated
  kill switch measures the wrong artifact: ``weekly binary downloads''
  can only mean \texttt{beadle-email} downloads, since that's the only
  binary distributed, and a download of the email MCP server says nothing
  about whether the signed-pipeline-daemon vision this document is named
  for has any pull. The 6-month-post-release clock lands 2026-09-13 ---
  see \faqref{faq:pnl} and \faqref{faq:metrics} for the current numbers
  and the recommendation to re-baseline it. \\
\bottomrule
\end{tabularx}
\renewcommand{\arraystretch}{1.0}

% ══════════════════════════════════════════════════════════════════════════════
% FEATURE APPENDIX
% ══════════════════════════════════════════════════════════════════════════════

\newpage
\section*{Feature Appendix}

Defines the scope boundary for the product. Every feature is explicitly
categorized to prevent scope creep and force prioritization decisions
upfront. This revision deliberately splits the usual Must/Should/Won't Do
into five categories instead of three --- Shipped, Built-Not-Distributed,
and Must Do stand in for what would otherwise be a single undifferentiated
Must Do, because the distribution gap between ``exists in the repo'' and
``a user can get it'' is this document's central finding (see the scope
note at the top).

\subsection*{Shipped (distributed to users)}

Features a user gets by installing Beadle today. This is the whole shipped
product --- \texttt{beadle-email} --- and nothing else in this appendix
reaches a user without it.

\begin{enumerate}[nosep,leftmargin=2.5em]
  \featureitem{MCP email server}{21 always-on tools plus 2 poll tools
    gated on config --- list, search, read, send,
    reply, mark, move (single + batch), attachments, verify signatures,
    MIME inspection, trust classification, address book, identity
    switching, inbox polling}\label{feat:mcp}
  \featureitem{Four-level transport trust}{trusted (Proton E2E), verified
    (PGP), untrusted (bad PGP), unverified (none). Inline PGP verification
    during message listing}\label{feat:trust}
  \featureitem{Per-contact rwx permissions}{read, write, execute per contact
    per identity}\label{feat:permissions}
  \featureitem{Multi-identity via ethos}{identity resolved per-request from
    ethos sidecar, with a plain-file fallback that needs no ethos install
    at all. Repo-local config, global fallback, mid-session
    switching}\label{feat:identity}
  \featureitem{Credential isolation}{secrets resolved at runtime from OS
    keychain (macOS Keychain, Linux pass/libsecret), secret files, or env
    vars. Never stored in config}\label{feat:credentials}
\end{enumerate}

\subsection*{Built, Not Distributed}

Real, tested code that no installed user can reach --- \texttt{make dist}
and \texttt{install.sh} build and ship \texttt{beadle-email} only.
Reaching users is gated on the Must Do items below, not on more building.

\begin{enumerate}[nosep,leftmargin=2.5em]
  \featureitem{Pipeline v1}{background daemon with IMAP poller, trust gate,
    command YAML loader with typed args, regex planner, sequential executor,
    auto-reply as terminal stage, crash-safe pipeline persistence, worker
    spawner via \texttt{claude -p --bare} with adversarial
    isolation}\label{feat:pipeline-v1}
  \featureitem{Email trigger via Proton Bridge}{receive instructions by email,
    verify sender identity, execute pipeline, return results by
    email}\label{feat:email}
  \featureitem{CLI runner}{execute tools (biff, jq, git, ethos) directly via
    \texttt{exec.Command} in milliseconds. Whitelist-only binary resolution.
    No shell, no LLM overhead for deterministic
    operations}\label{feat:cli-runner}
  \featureitem{Process/passthrough modes}{commands declare whether they
    transform the pipe (process) or read it as a side-effect (passthrough).
    Notifications don't destroy pipeline data}\label{feat:passthrough}
  \featureitem{Compound CLI steps}{chain binaries within a single command
    (jq | biff) via goroutine-per-step with io.Pipe. No shell
    invoked}\label{feat:compound}
  \featureitem{JSON pipe with schema validation}{structured data flows between
    stages. Commands declare output schemas. Daemon validates JSON between
    process-mode stages}\label{feat:schemas}
  \featureitem{Runner-conditional validation}{prompt and budget required for
    Claude commands only. Binary and steps for CLI only. Typed args for
    both}\label{feat:runner-validation}
\end{enumerate}

\subsection*{Must Do}

What has to be true before this is a product a user can evaluate, not
just code that passes tests. Launch-blocking.

\begin{enumerate}[nosep,leftmargin=2.5em]
  \featureitem{GPG command signing enforcement}{verify GPG signatures on
    command YAML files at daemon startup. Reject unsigned or tampered
    files. The verifier is implemented
    ({\small\texttt{VerifySignature}}, DES-034) and wired into the
    daemon's command loader (DES-035) --- this is the mechanism the
    document's headline is named for. An operator sets
    {\small\texttt{owner\_handle}} or {\small\texttt{owner\_gpg\_key\_id}}
    in a new {\small\texttt{daemon.json}} to name who may sign command
    files; a daemon that has not configured either loads zero commands at
    all, never unsigned ones --- no signing tool exists yet for a new user
    to sign their own command files with, so a fresh deployment runs no
    daemon-triggered commands until one exists}\label{feat:signing}
  \featureitem{\texttt{beadle init} onboarding}{abstract identity and
    signing setup so a new user never touches raw \texttt{gpg}. The
    Usability risk in this document is rated High specifically because
    this doesn't exist}\label{feat:init}
  \featureitem{Daemon distribution}{build and ship \texttt{beadle-daemon}
    the way \texttt{beadle-email} already ships --- cross-compiled
    binaries, an install path, a way to run it as a service. Without this,
    Pipeline v1 and v2 are permanently unreachable by anyone but the
    author}\label{feat:daemon-dist}
\end{enumerate}

\subsection*{Should Do}

Features for after the Must Do items ship. Not launch-blocking.

\begin{enumerate}[nosep,leftmargin=2.5em]
  \featureitem{LLM planner}{replace regex planner with Claude decomposition
    of natural language instructions into command
    sequences}\label{feat:llm-planner}
  \featureitem{Bridge health monitoring}{configurable health checks, a
    status command, outage logging, and a retry queue for pending result
    deliveries --- see \faqref{faq:tech-risks}}\label{feat:bridge-health}
  \featureitem{Service-manager integration}{a macOS LaunchAgent and a
    systemd unit that restart the daemon on
    failure}\label{feat:service-manager}
  \featureitem{Internal scheduler}{cron-like scheduled pipeline execution with
    no catch-up on missed runs}\label{feat:scheduler}
  \featureitem{Multi-channel triggers}{biff chat, webhooks as trigger channels
    alongside email}\label{feat:multi-channel}
  \featureitem{Key rotation management}{guided key rotation with expiry
    enforcement}\label{feat:key-rotation}
  \featureitem{Multi-provider email}{Fastmail, Gmail, self-hosted Dovecot.
    Standard IMAP/SMTP, no architectural lock-in}\label{feat:multi-provider}
\end{enumerate}

\subsection*{Won't Do}

Features explicitly excluded. Naming what you won't build prevents scope
creep and clarifies the product's identity.

\begin{enumerate}[nosep,leftmargin=2.5em]
  \featureitem{Multi-tenancy}{Beadle is a personal agent for one owner ---
    not to be confused with \featureref{feat:identity}'s multiple
    \textit{identities for that one owner}, which is shipped. Serving
    multiple distinct owners from one daemon adds complexity that
    conflicts with the single-owner trust model}\label{feat:no-multi-tenant}
  \featureitem{Signal or other messaging integrations}{email is the remote
    interface. Adding messaging platforms fragments the authentication model
    and increases attack surface}\label{feat:no-signal}
  \featureitem{Web UI}{a browser interface contradicts the local-first,
    CLI-native design philosophy and adds a large attack
    surface}\label{feat:no-web-ui}
  \featureitem{Cloud hosting or SaaS model}{Beadle runs on the owner's machine.
    A hosted version would undermine the core value proposition of data
    sovereignty and local control}\label{feat:no-cloud}
  \featureitem{Scheduler catch-up}{missed scheduled runs are skipped, not
    replayed. Catch-up logic adds complexity and ambiguity about which version
    of a command should execute}\label{feat:no-catchup}
  \featureitem{Expired or non-expiring command-signing keys}{all
    command-signing keys must carry an expiration date that has not yet
    passed. This is a security invariant, not a convenience
    trade-off}\label{feat:no-nonexpiring}
\end{enumerate}

% ══════════════════════════════════════════════════════════════════════════════
% REFERENCES
% ══════════════════════════════════════════════════════════════════════════════

\newpage
\printbibliography[title={References}]

\end{document}
